\section{Introduction}\label{sec:introduction}
A service agreement can return to the same public operating condition after different renewal reviews. Those histories need not permit the same future service: a customer who could have exited at an earlier review may have accepted a different continuation obligation. A provider that retains only the current public condition can therefore lose value even with Markovian demand, fixed prices, and no private information. The operational questions are how to represent the complete accepted continuation response, how much history an optimal protocol must retain, and whether randomization can compensate for a restricted record.

We study positive additive payments, separable strictly concave operating rewards, interval controls, and an exogenous recombining public graph. The single-quadratic case exposes the structure; a further theorem permits changing curvature and downward jumps in marginal reward through continuous piecewise-quadratic functions. There is one scalar commitment, and the exact compiler excludes intertemporal switching and coupled vector promises. Participation caps reflect a continuation price at public-vertex-specific barriers. A backward construction computes the complete parameterized response before the root promise is selected, without first obtaining an extensive-form optimum.

The optimization lineage is explicit. Probability weighting converts the unfolding into classical separable laminar allocation, as in \citet{Mjelde1983} and \citet{Tang1990}. Normalized memoization yields the same public-graph recurrence. Parametric piecewise-quadratic flow and polymatroid sensitivity are established subjects \citep{KlimmWarode2021,HarksKlimmPeis2018}. We give an explicit convex-cost-flow representation of the unfolding and distinguish its primal supply parameter from the dual price used by the compiler. Neither scalar marginal allocation, generic parametricity, nor memoization is claimed as a new principle. The contribution is a quantitative event budget measured in the compact public input, together with exact rational preprocessing and representation choices for that complete response.

The quadratic specialization has a linear global knot universe and a matching quadratic number of pieces across separately materialized vertex tables. A chain attains both orders. These are explicit-output bounds, not lower bounds for every possible shared representation. A shared arithmetic circuit retains local responses and graph links instead of all vertex tables, using linear space in the piecewise input after compilation at the cost of slower queries. The piecewise-quadratic extension bounds events by the local piece count plus the number of public caps. It covers genuine marginal jumps, retains exact reflection and rational certificates, and does not enlarge the alphabet of carried cap prices. Thus preprocessing space, persistent read-only representation, and execution memory are distinct resources.

Execution is governed by a finite deterministic output transducer driven by public transitions. Backward state minimization is classical \citep{BeanBirgeSmith1987,GivanDeanGreig2003,Peyriere2023}. The contract-specific conclusions are finite reachability from continuous controls, ordered behavioral classes, and a worst-case linear additional writable alphabet under a stated public-observation interface. We report that alphabet separately from combined public-state/class pairs and read-only coefficients. Distinct numerical prices can produce identical future behavior and therefore share a symbol. Conversely, a recombining renewal family requires a distinct symbol for each distinct optimal terminal tier. Finite policy graphs already represent dynamic contracts \citep{Zhang2012}; here the graph is exact for a fixed promise in a finite-horizon public-information model, not an approximation for adverse selection.

The renewal family provides both a deterministic memory-value frontier and a randomized extension. The deterministic optimum has contiguous cap-ordered cells under heterogeneous probabilities and costs. For a fixed randomized terminal codebook, adjacent lotteries give an exact concave-envelope formula. An exactly solved three-branch example shows that two symbols can have strictly smaller loss with randomization, even though three symbols remain necessary for the exact optimum. This result applies when participation is evaluated in conditional expectation before the private draw. The corresponding lottery is not asserted to satisfy every private realization's cap. The retained comparative static concerns ordered mean-preserving radial spreads, rather than arbitrary convex-order comparisons.

The randomized design problem can also be solved globally for quadratic terminal rewards and heterogeneous quadratic intermediate costs. We prove that an optimal codebook has cap-anchored levels except for its highest level, which may be genuinely continuous. An ordered dynamic program with an analytically optimized terminal interval then computes the complete restricted-budget frontier in polynomial exact arithmetic. A heterogeneous example demonstrates that constraining all levels to branch caps misses the optimum. Under participation imposed after every private draw, a separate theorem shows that randomization cannot improve the deterministic frontier. These results turn the distinction between expected and realized participation into a directly computable service-design choice.

The complete frontier admits a further structural acceleration. Its interpolation costs and terminal costs have ordered minimizers, even after the highest service level is optimized continuously. This permits both participation frontiers to be computed by monotone search rather than quadratic scans. A symbolic completion also preserves ordered minima in every submatrix, enabling an implemented classical linear-search algorithm for each budget layer. The search principle is classical; the new argument establishes its hypotheses for heterogeneous contractual costs and for the minimized terminal array. An exact same-input comparison isolates this improvement, preserves off-cap optimal service levels, and measures oracle work separately from machine-dependent runtime.

Computation checks the claimed parametric object rather than relying on an expensive unfolding as a novelty argument. The retained laminar, convex, public-promise, and repeated-query studies remain available. New exact tests cover changing curvature, marginal jumps, shared-circuit evaluation, root inversion, and bounded-conjugate certificates. An independently coded exhaustive active-set method constructs whole parametric paths on small unfolded quadratic problems and compares every affine region. This reference is deliberately small and is not a timing benchmark of a published specialized solver. The results provide reproducible mathematical and implementation evidence, not calibrated service data or a claim of universal computational superiority.

The broader accepted-adaptation theory on switching, comparator-adjusted rents, continuous states, and certified deployment is retained with its original assumptions and evidence. The present article adds exact closure, a storage--query tradeoff, and a contractual role for restricted-memory lotteries while preserving the earlier quadratic, deterministic, and shared-table results as identifiable parts of the theory.
